\subsection{Trace-dual coordinates and affine verification of kernel corrections}\label{sec:binary:coordinates}
All vector indices in this subsection belong to $I_\ell$ and all
field elements belong to $L=\F_{2^n}$, $n=2^\ell$.
The characteristic-two coordinate construction is related
to binary-field ring switching~\cite{DiamondPosen2024}
and to its masked versions~\cite{Kochergin2026,FrigoShelat2026}.

\begin{lemma}[Existence of an element of full binary degree]\label{lem:binary:full-degree-element}
There exists $\alpha\in L=\F_{2^n}$ such that
$[\F_2(\alpha):\F_2]=n$.
\end{lemma}
\begin{proof}
For any $a\in L$, the tower law applied to
$\F_2\subseteq\F_2(a)\subseteq L$
shows that $d:=[\F_2(a):\F_2]$ divides $n$.
If every $a$ had $d<n$, then, because $n=2^\ell$,
every such $d$ would divide $n/2$.
The field $\F_2(a)$ has exactly $2^d$ elements.
Repeating the product-permutation proof of
Lemma~\ref{lem:binary:frob-period} in this
finite field gives $a^{2^d}=a$.
Iterating Frobenius then yields
$a^{2^{n/2}}=a$ for every $a\in L$.
This would make all $2^n$ elements of $L$ roots of
$X^{2^{n/2}}+X$, a nonzero polynomial of degree
$2^{n/2}<2^n$. It contradicts the root-count lemma,
Lemma~\ref{lem:history:rootcount}.
Therefore some $\alpha$ has degree $n$.
\end{proof}

\begin{definition}[Evaluation basis]\label{def:binary:eval-basis}
Choose $\alpha\in L$ of degree $n$ over $\F_2$,
whose existence is guaranteed by
Lemma~\ref{lem:binary:full-degree-element}, and set
\[
 \beta_u:=K_u(\alpha),\qquad u\in I_\ell.
\]
\end{definition}

\begin{lemma}[Kernel values form a binary field basis]\label{lem:binary:eval-basis}
The family $(\beta_u)_{u\in I_\ell}$ is an $\F_2$-basis of $L$.
\end{lemma}
\begin{proof}
The polynomials $(K_u)_u$ form an $\F_2$-basis of
$\F_2[X]_{<n}$ by Proposition~\ref{prop:intro:basis}.
The evaluation map $P\mapsto P(\alpha)$ is injective on
that space: if a nonzero polynomial of degree $<n$
vanished at $\alpha$, it would contradict the minimality
of the degree-$n$ minimal polynomial.
Both domain and codomain have dimension $n$ over $\F_2$.
Thus the map is an isomorphism, and the images of the
kernel basis form a basis of $L$.
\end{proof}

\begin{definition}[Binary coordinate maps]\label{def:binary:coords}
For each $u\in I_\ell$, define the unique
$\F_2$-linear coordinate map $d_u:L\to\F_2$ by
\[
 x=\sum_{u\in I_\ell}\beta_u d_u(x)
 \quad\text{for every }x\in L.
\]
\end{definition}

\begin{lemma}[Coordinate extraction]\label{lem:binary:extract}
For $u,v\in I_\ell$, $d_u(\beta_v)=1$ if $u=v$ and $0$
otherwise; each $d_u$ is additive.
\end{lemma}
\begin{proof}
The kernel-basis expansion of $\beta_v$ has coefficient $1$
at index $v$ and $0$ elsewhere; uniqueness of coordinates
gives the first assertion. If $x=\sum_v\beta_v a_v$ and
$y=\sum_v\beta_v b_v$, with $a_v,b_v\in\F_2$, then
$x+y=\sum_v\beta_v(a_v+b_v)$. Uniqueness gives
$d_u(x+y)=a_u+b_u=d_u(x)+d_u(y)$.
\end{proof}

\begin{lemma}[Trace pairing is nondegenerate]\label{lem:binary:trace-pairing}
For every nonzero $a\in L$ there exists $x\in L$ with
$\Tr(ax)=1$.
\end{lemma}
\begin{proof}
Choose $y\in L$ with $\Tr(y)=1$, using
Lemma~\ref{lem:binary:trace-surj}. As $a\neq0$,
$a^{-1}$ exists. Put $x=a^{-1}y$; then
$\Tr(ax)=\Tr(y)=1$.
\end{proof}

\begin{theorem}[Existence and uniqueness of trace-dual coordinates]\label{thm:binary:dual-basis}
There is a unique family $(\eta_u)_{u\in I_\ell}$ in $L$ such that
\[
 \boxed{\Tr(\eta_u\beta_v)=\delta_{uv}.}
\]
For every $x\in L$,
\[
 \boxed{d_u(x)=\Tr(\eta_u x).}
\]
\end{theorem}
\begin{proof}
Define the $\F_2$-linear map
$\Phi:L\to\operatorname{Hom}_{\F_2}(L,\F_2)$ by
$\Phi(a)(x)=\Tr(ax)$. Linearity follows from
Lemma~\ref{lem:binary:trace-add}.
Lemma~\ref{lem:binary:trace-pairing} proves that $\Phi$
is injective. Its domain has dimension $n$ over $\F_2$;
its codomain also has dimension $n$, since a linear functional
is uniquely determined by its values on the $n$ basis elements
$\beta_v$. Hence $\Phi$ is bijective.
For each coordinate functional $d_u$, choose the unique
$\eta_u=\Phi^{-1}(d_u)$. Then
$\Tr(\eta_u\beta_v)=d_u(\beta_v)=\delta_{uv}$
by Lemma~\ref{lem:binary:extract}, and
$\Tr(\eta_u x)=d_u(x)$ for all $x$ by construction.
Uniqueness of each $\eta_u$ follows from injectivity
of $\Phi$.
\end{proof}

\begin{definition}[Binary transpose and kernel correction]\label{def:binary:transpose}
For $p=(p_v)_{v\in I_\ell}\in L^{I_\ell}$ define
\[
 c^p_u:=\sum_{v\in I_\ell}\beta_vd_u(p_v),\qquad
 C_p(X):=\sum_{u\in I_\ell}c^p_uK_u(X).
\]
For $\theta\in L^\ell$ define
\[
 \Delta_p(\theta):=\sum_{u\in I_\ell}
       \lambda_u(\theta)c^p_u
       =D_{C_p}(\theta).
\]
\end{definition}

\begin{lemma}[Trace expression for the correction]\label{lem:binary:delta-trace}
For every $p\in L^{I_\ell}$ and $\theta\in L^\ell$,
\[
 \boxed{\Delta_p(\theta)=
  \sum_{u,v\in I_\ell}
   \lambda_u(\theta)\beta_v\Tr(\eta_u p_v).}
\]
\end{lemma}
\begin{proof}
Expand $\Delta_p(\theta)$ using
Definition~\ref{def:binary:transpose}:
\[
 \Delta_p(\theta)=
 \sum_u\lambda_u(\theta)\sum_v\beta_vd_u(p_v).
\]
Interchange the two finite sums and replace
$d_u(p_v)$ by $\Tr(\eta_up_v)$ using
Theorem~\ref{thm:binary:dual-basis}.
\end{proof}

\begin{definition}[Public affine coefficients]\label{def:binary:affine-H}
For $\theta\in L^\ell$ and $u,v,k\in I_\ell$ put
\[
 H_{v,k}(\theta):=\sum_{u\in I_\ell}
       d_k\bigl(\beta_v\lambda_u(\theta)\bigr)\eta_u.
\]
For a candidate coefficient table $c:I_\ell\to L$ define
\[
 D_c(\theta):=\sum_u\lambda_u(\theta)c_u
\]
and
\[
 V_k(p,c,\theta):=\eta_kD_c(\theta)
      +\sum_{v\in I_\ell}H_{v,k}(\theta)p_v.
\]
The coefficients $H_{v,k}(\theta)$ depend only on public
parameters and $\theta$, not on the hidden pads $p_v$.
\end{definition}

\begin{lemma}[Affine transpose trace identity]\label{lem:binary:affine-trace}
For every $p,\theta,k$,
\[
 \Tr\left(\sum_v H_{v,k}(\theta)p_v\right)
   =d_k\bigl(\Delta_p(\theta)\bigr).
\]
\end{lemma}
\begin{proof}
Expand the definition of $H$ and use
$\F_2$-linearity of trace:
\[
 \Tr\left(\sum_vH_{v,k}p_v\right)
 =\sum_{u,v}
    d_k(\beta_v\lambda_u(\theta))\,\Tr(\eta_up_v).
\]
Each factor $\Tr(\eta_up_v)$ lies in $\F_2$.
Using additivity and $\F_2$-linearity of $d_k$,
the last sum equals
\[
 d_k\left(\sum_{u,v}
   \beta_v\lambda_u(\theta)\Tr(\eta_up_v)\right).
\]
The field element inside $d_k$ equals
$\Delta_p(\theta)$ by Lemma~\ref{lem:binary:delta-trace}.
\end{proof}

\begin{lemma}[Trace coordinates of affine residuals]\label{lem:binary:V-trace}
For all $p,c,\theta$ and $k$,
\[
 \boxed{\Tr(V_k(p,c,\theta))
      =d_k\bigl(D_c(\theta)+\Delta_p(\theta)\bigr).}
\]
\end{lemma}
\begin{proof}
By the definition of $V_k$ and trace additivity,
\[
 \Tr(V_k)=\Tr(\eta_k D_c(\theta))
  +\Tr\left(\sum_vH_{v,k}p_v\right).
\]
The first term equals $d_k(D_c(\theta))$ by
Theorem~\ref{thm:binary:dual-basis}; the second equals
$d_k(\Delta_p(\theta))$ by
Lemma~\ref{lem:binary:affine-trace}.
Additivity of $d_k$ combines them into the claimed formula.
\end{proof}

\begin{theorem}[Exact affine trace consistency]\label{thm:binary:affine-consistency}
For fixed $p,c,\theta$,
\[
 \boxed{
 D_c(\theta)=\Delta_p(\theta)
 \ \Longleftrightarrow\
 \Tr(V_k(p,c,\theta))=0\quad\forall k\in I_\ell.
 }
\]
\end{theorem}
\begin{proof}
In characteristic two the equality
$D_c(\theta)=\Delta_p(\theta)$ is equivalent to
$D_c(\theta)+\Delta_p(\theta)=0$.
Lemma~\ref{lem:binary:V-trace} identifies
$\Tr(V_k)$ with the coordinate
$d_k(D_c+\Delta_p)$. A field element is zero exactly
when every coordinate in the basis $(\beta_k)$ is zero,
by Definition~\ref{def:binary:coords}. This proves both directions.
\end{proof}

\begin{lemma}[Mismatch gives a nonzero multilinear polynomial]\label{lem:binary:nonzero-mismatch}
If $c\neq c^p$ as tables, then
\[
 E_{c,p}(\Theta):=D_c(\Theta)+\Delta_p(\Theta)
   =\sum_{u\in I_\ell}(c_u+c^p_u)\lambda_u(\Theta)
\]
is a nonzero polynomial in $L[\Theta_0,\ldots,\Theta_{\ell-1}]$
of total degree at most $\ell$.
\end{lemma}
\begin{proof}
Each $\lambda_u$ is a product of $\ell$ affine factors,
giving the degree bound. Choose $u_0$ with
$c_{u_0}\neq c^p_{u_0}$.
Evaluate at the Boolean point $\Theta=u_0$. By
Lemma~\ref{lem:intro:interpolation} and
Lemma~\ref{lem:binary:eq-weights},
$E_{c,p}(u_0)=c_{u_0}+c^p_{u_0}\neq0$.
Hence the polynomial is nonzero.
\end{proof}

\begin{corollary}[Random-challenge detection of a coefficient mismatch]\label{cor:binary:mismatch-bound}
If $c\neq c^p$ is fixed before uniform
$\theta\in L^\ell$ is sampled, then
\[
 \Pr_\theta[D_c(\theta)=\Delta_p(\theta)]
 \leq \min\{1,\ell/q\},\qquad q=2^n.
\]
\end{corollary}
\begin{proof}
By Lemma~\ref{lem:binary:nonzero-mismatch},
$E_{c,p}$ is nonzero of total degree at most $\ell$.
The event on the left is $E_{c,p}(\theta)=0$.
Apply the multivariate Schwartz--Zippel bound proved in
Lemma~\ref{lem:history:schwartz-zippel}.
The fixed-before-challenge hypothesis ensures the
polynomial is independent of the uniform $\theta$.
\end{proof}

\begin{remark}[Meaning of this part of the construction]
The coefficients $H_{v,k}(\theta)$ turn the
$\F_2$-linear coordinate relation into \emph{affine equations
over the extension field}. This is an algebraic ingredient
of our proposed kernel-coordinate checker; it does not
by itself establish hiding of a committed witness or
the soundness of the complete ring-switching protocol.
\end{remark}
